\chapter[Why the Reformation Divided Europe]{Why Did the Reformation Tear Europe Apart?}
\section{An Indulgence Certificate}
Pick up another object. This time it is neither a stone tool nor a coin, but a piece of paper.

It is about the size of your palm, with printed ornament and coats of arms around the edges. The body is in Latin, with several blank lines for someone to enter a name and date in ink. Five hundred years ago, papers like this were sold individually in German marketplaces. Buyers received written certification: because they had paid the Church this sum and performed the prescribed repentance, their punishment in purgatory---or that of a deceased relative they named---could be reduced by a certain number of years.

First we need to explain ``purgatory''. In the medieval European Catholic picture of the world, heaven and hell were not the only destinations after death. Most ordinary people were neither holy enough to ascend straight to heaven nor wicked enough to fall into hell for ever. They went to purgatory, where fire-like punishment gradually expiated the sins left unsettled during life. Only then could they enter heaven. This might take hundreds or thousands of years. The living could help souls there obtain a ``reduced sentence'' through prayer, good works and Masses. Within this system, an indulgence was the form of reduction that money could buy.

Weigh that paper in your hand. Almost weightless, it claimed to shorten punishment measured in centuries in another world. It was a receipt: the money went to Rome, where St Peter's Basilica, then the world's most magnificent church, was under construction. It was also a contract. One party was the ordinary person paying; the other regarded itself as God's general representative on earth: the Roman Church.

Inside this paper lies our question: why did arguments over it split Europe in two within little more than thirty years, set neighbours killing neighbours, put fathers and sons in opposing camps, and draw millions into successive wars fought in God's name?

Paper alone cannot exert such force. The force came from other things. Following this sheet, we shall uncover them one by one.

\section{A Market Square in 1517}
Come into a scene with me. First, a clarification: the following reconstruction draws on surviving sermons, accounts and letters. I have invented the particular names, but every element of the scene has documentary support.

The time: around 1517, in early spring, as the snow thaws. The place: the market square of a small German town. These German lands were not yet ``Germany'', but a patchwork of hundreds of principalities, free cities and ecclesiastical territories, nominally sharing allegiance to the loose framework called the Holy Roman Empire.

A temporary preaching platform stands in the square, with a cross before it and the pope's coat of arms hanging above. A mendicant friar selling indulgences speaks from it in a booming voice. Such travelling sales campaigns were an established business, with budgets, divided responsibilities and commissions. Beside the platform is a table, with an accounts clerk behind it, quite possibly employed by the Fugger family, the great Augsburg bankers. The division of the proceeds was already agreed: half would go to Rome for St Peter's; the other half would stay locally to repay the enormous loan that Albrecht, Archbishop of Mainz, had obtained from the Fuggers to purchase his office. From the outset, this ``spiritual'' transaction was attached to a thoroughly real financial chain.

The friar is preaching about purgatory. He offers stories, not abstractions. Your mother died three years ago: have you considered where she is now? In the flames. Her sins remain unexpiated; the fire burns her. You sleep soundly every night while she waits in the fire. For what? For you to do something entirely within your power. The friar raises the printed sheet. According to contemporary reports, salesmen of this kind used a catchy slogan, roughly: the moment a coin rings in the chest, a soul flies out of purgatory. Whether any particular salesman actually uttered those exact words is now difficult to establish. Nevertheless, they capture the sales pitch: simple, direct and accompanied by an audible chink.

A peasant woman stands in the queue, flour on her apron and a few coins clenched in her hand. She saved them by selling eggs, intending to buy cloth for her younger daughter's winter clothes. Her late husband liked drinking and became abusive when drunk. According to the priest, a soul like his might spend an incalculable number of years in purgatory. She hands over the coins, gives his name and watches the clerk carefully write it down. Her hands tremble as she takes the paper. Has she been deceived? In terms of her own feelings at that moment, probably not: she has bought something she sincerely believes works, and feels she can finally do something for her husband. Remember this detail. The indulgence business depended not merely on sellers' persuasive words, but on buyers' genuine, burdensome love and fear.

A parish priest stands on the edge of the crowd, looking grim. His concern is money, not doctrine: every additional copper entering the travelling seller's chest means one fewer for his own church's collection box. Hundreds of kilometres away in Wittenberg, a monk and theology professor named Martin Luther has precisely the opposite concern. He worries not about the money, but that this trade is turning repentance, a matter of the soul, into a business where one can receive change.

Remember this square: the platform, the money chest, the peasant woman, the bank clerk and the grim-faced priest. The earthquake that later tore Europe apart had its epicentre beneath ordinary marketplaces like this one.

\section{Ninety-Five Questions on a Door}
Now we can set out the conflict.

On 31 October 1517, Martin Luther published his Ninety-Five Theses in Wittenberg. The classic image, handed down for five centuries, shows him swinging a hammer to nail the document to the Castle Church door, which served rather like a university noticeboard. Honesty requires a qualification: some modern historians question whether the details of the nailing scene are entirely authentic, since it rests principally on later accounts. What is undisputed is that Luther wrote the ninety-five propositions and sent them to the Archbishop of Mainz, who shared in the indulgence revenue.

Another clarification: the Ninety-Five Theses were written in Latin as a public invitation to an academic disputation. Their tone was relatively restrained for the time. Luther did not yet intend to ``rebel'', much less found a new Church. He wanted to settle the question of indulgences within the existing Church.

Yet the questions he uncovered went much deeper. Beneath the theses lay three progressively more fundamental layers:

\textbf{First: what exactly could the pope remit?} According to the Church's own theory, indulgences concerned the remission it could grant of ``temporal punishment''. Salesmen, however, presented them as a universal remedy: they could rescue souls in purgatory and absolve sins great and small. Luther asked which promises were true: those in official Church documents or those in the sales pitch? If even the Church's documents did not support this way of selling, why was its entire machinery protecting the campaign?

\textbf{Second: how can a person be accepted by God?} Here was the real explosive. Years of spiritual struggle had led Luther to conclude that people did not purchase God's acceptance through good works, indulgences or accumulated merit. If good deeds were payment to God, their relationship became a commercial exchange. Would those too poor to afford the purchase then be excluded for ever? He argued for ``justification by faith'': people were counted righteous through faith in God's grace, not through points accumulated by their actions. If this was right, indulgences were more than a disputed business: their entire direction was mistaken. They sold something that was not the Church's to control.

\textbf{Third: who was entitled to settle the first two questions?} The Church answered: the Church, naturally---the pope and general councils, drawing on fifteen centuries of interpretative tradition. In the controversies of the following years, Luther was pushed step by step into a corner. His eventual answer grew increasingly radical: Scripture itself was the final standard, and every believer could, and should, read it personally.

Notice the structure. Each question goes deeper than the last. The first says, ``Your business is keeping fraudulent accounts'', something many within the Church acknowledged. The second says, ``Your entire points system is misdirected'', shaking centuries of theology. The third says, ``You have no right to monopolise the answers'', shaking the whole structure of authority. Historians often debate the Reformation's ``real cause''. In fact, each layer had its own way of igniting events. Later we shall use a particular tool to separate them.

For now, leave the chapter's central question hanging: \textbf{how did one monk's doubts become the sound of cannon on a battlefield?} Critics of Church corruption had never been lacking over the preceding centuries. Most were suppressed or burned. Why did this challenge in 1517 keep burning more fiercely instead of dying down?

\section[Other Voices]{A Word for Indulgences, and Other Voices}
Before proceeding, let us repeat an exercise used throughout this book: give the ``losers'' their full argument. Indulgences and the Church that sold them have become almost synonymous with corruption in later accounts. But if you can only repeat the victors' verdict, you have not really understood the history.

\textbf{First, make the strongest possible case for indulgences.}

First, the system responded to real fears; it did not manufacture them from nothing. Medieval people genuinely believed in purgatory and worried about their dead relatives. Within that faith, indulgences gave ordinary people an achievable course of action. You need not practise lifelong austerity like a monk: paying for a Mass or making a donation could help someone you loved. To call this simply a swindle amounts to calling entire generations fools. They were not fools. They inhabited a picture of the world they believed to be true.

Second, it had an internally coherent theology. Officially, indulgences never meant ``buying your way into heaven''. They presupposed sincere repentance and confession, and remitted only the ``temporal punishment'' still requiring satisfaction afterwards. The Church also taught a ``treasury of merit'': Christ and the saints had accumulated more merit than they needed, and the Church, as its steward, could redistribute it to believers who lacked sufficient merit. This theory was internally consistent. You may reject its premises, but cannot call it devoid of logic.

Third, the money did not all enter private pockets. Building St Peter's was regarded as a public undertaking for Christendom, much as a city today might solicit donations for a cathedral or major museum. Dismissing everything as ``Church corruption'' is too easy.

Fourth, the Church clearly issued a warning that events partially vindicated: allowing everyone to interpret Scripture independently would produce disorder. This was not pure intimidation. The Reformation did bring radical groups such as the Anabaptists, demanding a fresh start; peasant wars fought under religious banners; and chaotic disputes in which rival voices declared one another heretical. Looking forward from 1517, the Church's claim that a unified interpretation sustained order had a basis in reality. It chose, however, to maintain that unity through suppression rather than argument.

\textbf{Now hear the reformers.}

Luther's central claim, in his own terms, was that Scripture had ``bound'' his conscience. At the Diet of Worms in 1521, he was ordered to retract his writings. According to the traditional account, his answer was substantially this: unless refuted by Scripture or clear reason, his conscience remained captive to God's word; he could not and would not recant. Historians have reservations about the exact authenticity of the famous ``Here I stand; I can do no other''. Yet it captures the scene's spirit: an individual claimed that something outweighed emperor, pope and personal safety---the text he had read himself, and his own conscience. This was a strikingly new stance at the time.

\textbf{There were third and fourth voices too, often drowned out by theologians and reformers.}

One belonged to the German princes. To them, the Luther controversy was above all an opportunity. Adopting Protestantism could legitimise confiscating extensive Church lands, retaining taxes otherwise sent to Rome and escaping obedience to a foreign power influenced by the pope. This does not make every prince an opportunist: some were sincere converts. But everyone could hear the clatter of calculations of self-interest.

Ordinary people's voices were more complicated. In the German Peasants' War, which began in 1524, thousands carried flails and pitchforks and cited Scripture against serfdom's oppression. Their reasoning was plain: if the Gospel proclaimed equality, why were they their masters' property? Intriguingly, Luther, whom they claimed as an ally, sided with the princes and fiercely condemned the revolt, essentially demanding its suppression. About 100,000 peasants were killed in two years. Once ``the language of liberation'' enters reality, who may interpret it, and how far that interpretation may go, immediately becomes another struggle for power. Breaking the old Church's monopoly did not remove the question of interpretative authority; it moved the battlefield.

Can you hear the difference? This is not a story of ``bad people selling false certificates and good people exposing the scam''. It is an impasse in which each party holds part of an argument and pursues its own calculations. Understanding that impasse brings us closer to history than pronouncing judgement.

\section[Thinking Bridge: Three Strands]{Thinking Bridge: Three Strands Twisted into One Rope}
Here is a tool for analysing explosive historical changes: the Reformation, the French Revolution or the internet revolution.

It is simple: \textbf{separate the causes of a great change into three strands, then examine how they intertwine.}

\textbf{The first strand: ideas.} What changed in people's thinking? What were the new claims, and where did they clash with old ones? For the Reformation, these were theological propositions such as justification by faith and Scripture's supreme authority.

\textbf{The second strand: interests.} Who gained or lost what? How were money, land, taxes, power and status redistributed? Here the strand includes princely estates, Rome's taxes, bank loans and townspeople's tax burdens.

\textbf{The third strand: conditions.} What technologies or circumstances made something formerly containable impossible to suppress? Here the strand includes printing, urban networks, universities and slowly rising literacy.

The crucial move is not separation but intertwining: \textbf{looking at only one strand guarantees an elegant but mistaken explanation.}

Let us test an explanation with counterexamples. Someone says, ``Printing caused the Reformation.'' That sounds plausible. Luther's theses spread through German towns within weeks; pamphlets and woodcut cartoons flew like snowflakes, inconceivable in a manuscript world. But consider two counterexamples. First, movable type appeared earlier in the East. Korean metal type, Song Chinese movable type and earlier woodblock printing spread Buddhist scriptures widely, yet did not generate a ``Buddhist Reformation''. Second, about a century before Luther, the Czech Jan Hus proposed very similar Church reforms, only to be burned at the stake and have his movement suppressed. The ideas already existed; the conditions were missing. Printing was an important condition, but could neither supply a reform's content nor automatically secure its victory.

Conversely, ``because the Church was too corrupt'' is inadequate. The Church in 1500 was not an order of magnitude more corrupt than in 1350. Why did the explosion wait until 1517? Nor is ``because princes wanted Church land'' sufficient. Princes had wanted land for ages. Without the strand of ideas, they could not seize the banner of ``freedom of conscience'' or mobilise people in the marketplace sincerely anxious about their mothers' souls.

Only when twisted into one rope could these strands pull history forward: \textbf{ideas supplied reasons, interests supplied motivation, and conditions supplied possibility.} The next time you encounter an argument over an event's ``fundamental cause'', whether in a history question or a news comment section, ask whether the speaker is grasping only one strand.

\section{Reading a Few Original Lines}
Now read a short primary source: the document that ignited it all.

Written in 1517, the Ninety-Five Theses were originally Latin propositions for academic debate. Chinese translations vary slightly in wording. Here are the main points of several crucial theses, with their numbers so you can check them yourself later:

\begin{quote}
Thesis 27: Those who say that a soul flies from purgatory as soon as a coin rings in the chest are preaching inventions.
\end{quote}
\begin{quote}
Thesis 28: It is clear that a coin ringing in the chest can only increase greed and avarice.
\end{quote}
\begin{quote}
Thesis 36: Every truly repentant Christian can receive full remission of guilt and punishment, even without an indulgence.
\end{quote}
\begin{quote}
Thesis 82: Since the pope is now wealthier than the richest man, why does he not build St Peter's with his own money, instead of that of poor believers?
\end{quote}
Read these together and notice their increasing range. Theses 27 and 28 still attack salesmanship: even your own theology does not support souls ascending at the sound of coins. Thesis 36 goes deeper, saying remission rests on genuine repentance rather than certificates. Indulgences change from an ``over-advertised product'' to a ``fundamentally misguided product''. Thesis 82 strikes hardest, bypassing theology to ask about money: if you care so much about souls, why not pay yourself? It translates faith into finance, in terms every European taxpayer could understand.

Now consider the groundwork Luther laid. In arguing for justification by faith, he repeatedly cited the New Testament's Epistle to the Romans, especially a sentence in its first chapter:

\begin{quote}
``The righteous shall live by faith.'' (The Bible, Romans, chapter 1; translated from the Chinese Union Version.)
\end{quote}
Paul was quoting the older Jewish Book of Habakkuk. Notice the words: living by ``faith'', not by ``certificates'' or ``merit points''. Repeated reflection on this sentence led Luther to his revolutionary conclusion: God's acceptance depended on grace and faith, not services sold over a Church counter.

Why was this document more dangerous than countless attacks on the Church over earlier centuries? It combined two things others had not achieved together. It used the Church's own most revered texts to challenge its most profitable business, and every blow landed where ordinary people could understand: the claim that your mother awaited your payment in the flames was a scam. It lacked neither scholarship nor destructive force.

\section{One Earthquake, Three Explanations}
Return to our central question: how did one monk's doubts become cannon fire? Historians offer several explanations of why the Reformation began and succeeded. They are not interchangeable; each has evidence and blind spots. Let us compare three representative accounts. Remember our distinction between four kinds of content? What happened---theses, controversies and wars---belongs to the evidential level. Why it happened is where the following arguments clash.

\textbf{Explanation one: theology---a dispute over direction within a faith.}

Here the central issue is salvation: on what basis does God accept a person? One side answers faith together with good works, sacraments and ecclesiastical hierarchy; the other, faith alone, grace alone and Scripture alone. The evidence is strong: participants understood matters this way and accepted imprisonment, exile and death over these differences. Subsequent divisions between denominations largely followed theological disagreements: whether Christ's presence in the Eucharist was real or symbolic, whether baptism belonged to infants or adults. These were genuine questions. The limitation is time and place. Why did similar theological dissatisfaction fail in Hus's era but succeed in Luther's? Why did some princes fight for faith, others for \textbf{calculations about faith}? Theology alone cannot explain ``why here and now''.

\textbf{Explanation two: political economy---faith was the banner; redistribution was the substance.}

On this account, the Reformation fundamentally transferred wealth and power. Princes used Protestantism to confiscate Church estates, retain taxes and escape Rome's and the emperor's constraints. Townspeople escaped ecclesiastical exactions. Northern European monarchies subordinated churches to the state. There is plenty of evidence: Henry VIII launched England's Reformation over divorce and succession, with theological differences almost an afterthought. The map of German princely allegiances closely matched their interests. This sharp explanation illuminates ``why here and now'', but its blind spot is equally clear: participants who gained nothing and paid dearly. What did peasants killed in the revolt, Anabaptists at the stake or clergy abandoning comfort for exile expect to gain? Translating everything into interests resembles calling our peasant woman a ``brainwashed consumer''. It makes sense, but misses the real warmth of the coins in her palm.

\textbf{Explanation three: communication---an old fire amplified by technology.}

Here criticism of the Church had smouldered for centuries; the new variable was communication. Printing spread the theses across Germany within weeks. Woodcut cartoons let the illiterate ``read'' satire of the pope. Networks of towns and trade routes carried news faster than prohibitions. This addresses both ``why now''---printing had matured only decades earlier---and ``why here'': dense German towns, numerous princes and weak central authority made bans difficult to enforce. We have already tested its limitation: earlier, more extensive printing of Buddhist texts in the East produced no comparable explosion. Conditions amplify content; they do not replace it.

Each explanation grasps one strand: theological ideas, political-economic interests or communicative conditions. A mature view does not choose a camp, but recognises that without any one strand this fire would not have caught. Together they pulled Europe through half a century of upheaval. This is not evasive compromise. It is testable: it predicts that ideas without interests, as with isolated missionaries; interests without ideas, as with straightforward land annexation; and conditions without content, as in the East's printing of Buddhist scriptures, would not produce a Reformation. History bears this out.

\section{War, States and Everyday Life}
The seeds of reform produced not one uniform Protestantism, but a thicket of different plants.

In Zurich, Zwingli cleared sacred images from churches, pursuing a more thorough reform than Luther. In Geneva, Calvin wrote the Institutes of the Christian Religion and made ecclesiastical discipline part of daily civic life. His emphasis that God had already determined who would be saved---``predestination''---spread with his followers through France, the Netherlands and Scotland. The Anabaptists went further, advocating voluntary adult baptism and refusing allegiance to secular power. Catholics and Protestants alike suppressed them; many leaders were drowned, a deliberate executioner's mockery of those who ``rebaptised''. England took a fourth route: the king declared separation from Rome from above, while doctrinal change was smallest. Protestantism never had a single voice. Luther and Zwingli's bitter face-to-face dispute over the Eucharistic bread and wine provided an early public demonstration of the cost of ``everyone reading Scripture for themselves''.

The Catholic Church did not sit idle. Meeting intermittently for eighteen years from 1545, the Council of Trent fixed doctrinal boundaries while pursuing substantial reforms: prohibiting abuses in indulgence sales, requiring bishops to reside in their dioceses and establishing seminaries to train priests. The newly founded Jesuits opened schools and undertook overseas missions, winning back many hearts. Historians call this the ``Catholic Reformation'' or ``Counter-Reformation''. Notice the difference: ``counter'' casts Catholicism as responding defensively; ``reformation'' acknowledges internal demands for renewal that predated Luther and were not simply reactions to him.

Then came cannon fire. The German Peasants' War of 1524--1525 killed about 100,000 peasants. In the St Bartholomew's Day massacre in Paris in 1572, Catholics killed thousands of Protestant Huguenots, triggering decades of French religious civil war. The Thirty Years' War, 1618--1648, devastated the German lands. From a pre-war population of roughly sixteen to twenty million, millions died through warfare, famine and disease; some areas lost over a third of their inhabitants. Wars begun over doctrine often ended with mercenaries fighting for wages and great powers for supremacy. Late in the Thirty Years' War, Catholic France financed Protestant Sweden's attacks on the Catholic emperor. Theological banners faded before political calculations.

When fighting could no longer continue, new rules had to be negotiated. The Peace of Augsburg in 1555 established ``cuius regio, eius religio'': whose territory, his religion. Princes chose a confession; subjects followed or emigrated. The Peace of Westphalia, ending the Thirty Years' War in 1648, extended this arrangement across Europe: rulers governed their own territories without foreign interference. Today's international-law textbooks commonly trace the ``sovereign-state system'' to that year. Counterintuitively, the outlines of the modern state were largely forged by religious war---not because people became tolerant, but because neither side could destroy the other. They invented rules of mutual non-interference, demoting faith from ``a dispute over universal truth'' to ``domestic affairs''.

Finally came the easily overlooked, perhaps deepest change: discipline at the level of meals and schooling. Protestantism demanded that everyone read Scripture, which required literacy. Protestant territories expanded schools, and literacy slowly increased. After Trent, Catholics likewise strengthened parish schools, preaching and registration: baptisms, marriages and deaths all entered the records. Both sides competed to infuse daily life with religion through prayers before meals, Sunday worship, catechisms and records of home visits. Historians call this ``confessionalisation'': opposing denominations independently regulated believers more closely and uniformly, incidentally training disciplined, literate, registered subjects for emerging states. Max Weber later advanced a famous, disputed claim: Protestantism, especially Calvinism, treated secular work as a ``calling'' in God's service, unintentionally encouraging the spirit of modern capitalism. Scholars still defend and challenge this argument. Here it illustrates how daily discipline can reshape economic behaviour, rather than serving as a settled conclusion.

Notice four things we have deliberately kept apart. Wars and treaties concern what happened. Confessionalisation and the Weber thesis are explanations of why, each with supporters. Participants understood their wars as defending the true faith. Our judgement that the wars ``gave birth to the modern state'' rests on hindsight. Mix these together and history becomes a muddle.

\section[Further Challenge: Where Authority Lives]{Further Challenge: Where Does Authority Live?}
Now for the chapter's weightiest theoretical tool. Beneath every preceding dispute---indulgences, theses, the Eucharist and war---lies one question: \textbf{where does authority live?} Who decides, and on what grounds?

Max Weber distinguished three types of reason for obeying authority. Everyday examples help. People obey a hereditary chief because ``it has always been this way'': traditional authority. They follow an exceptionally inspiring figure because ``this person is extraordinary'': charismatic authority. You obey traffic police and examination rules not through tradition or personal devotion, but because you recognise the legitimacy of their rules and procedures: legal-rational authority.

Applied to the Reformation, this tool clarifies the picture. The medieval Roman Church combined traditional and legal-rational authority: fifteen centuries of apostolic succession, alongside an immense system of canon law and ecclesiastical hierarchy. Luther was a classic charismatic figure. Yet he sought to relocate authority from institutions to a text. Scripture alone: the text was supreme.

But here is the deeper challenge: \textbf{texts cannot speak for themselves; people must read and translate them.} As soon as authority moved into a book, new struggles began. Who could read? Who owned books? Into which language would they be translated, and using which words?

Translation became one of the Reformation's sharpest yet least visible battlefields. Sheltering at the Wartburg, Luther translated the New Testament into German, later completing the whole Bible. He deliberately used everyday rather than scholarly German so that ``weaving women and schoolchildren'' could read it. England's Tyndale had begun the same work earlier. He deliberately rendered certain Greek terms as ``elder'' and ``congregation'', rather than the traditional ``priest'' and ``church''. Every choice rewrote power relations: without the word ``priest'', the clerical hierarchy lost a little of its sacred aura. Tyndale was eventually burned, his translation at the heart of the charges. A translation could get someone killed: translation was never merely linguistic.

Look beyond Europe and authority's entanglement with translation becomes a human-wide question. After Buddhism reached China, generations of translators from the Eastern Han to the Tang argued over rendering Sanskrit. Xuanzang travelled west for scriptures, directed translation workshops and established the ``five kinds not to translate'': five categories better transliterated than forcibly translated, lest meaning be distorted. Another civilisation treated fidelity to an original as a matter of faith: different methods, similar anxieties. In Islamic tradition, the Arabic Qur'an holds supreme status. Many religious scholars regard translations only as interpretations of meaning, not the authentic scripture itself. The Islamic world therefore experienced almost no revolution in which a translation's authority displaced the original, and authority developed quite differently. Three traditions arranged originals and translations in three ways, producing different relations between religion and power. This is not a ranking of advancement. It reminds us that Luther made a historically particular choice, whose mixed consequences---dispersed interpretative authority---continue to unfold.

Weber's tool offers one last counterintuitive insight. Luther broke an old authority without producing a utopia of free interpretation. Protestant denominations soon developed their own doctrinal scrutiny and church discipline, sometimes regulating life more closely than Catholics. Authority was not abolished; it changed address, from Rome to princely courts, city councils and family Bibles. After the charismatic Luther died, his movement followed the familiar path of institutionalisation. Weber observed unsentimentally that this is almost every charismatic authority's fate. Test that prediction against any historical movement promising to shatter old authority.

\section{The Indulgence in Your Pocket}
This five-hundred-year-old question is in your pocket now.

Consider what printing did: bring the costs of producing and obtaining texts close to zero, breaking authority's information monopoly. The internet has done it again on an enormously greater scale. Then everyone could ``read Scripture themselves''; now everyone can ``look up papers'', ``inspect raw data'' and ``do their own research''. Authority must move once more. Textbooks, experts and institutional media no longer have the final say; ordinary people who can search and compare have tools for challenging them.

And then? The old scenes repeat. Dispersed interpretation corrects errors: patients consult research and challenge misdiagnoses; internet users compare original documents and expose inaccurate reporting. It also produces disorder: conspiracy theories use ``doing your own research'' as a shield; pseudoscience opens with ``authoritative institutions are all lying to you'', much as radical sects began with ``Rome is lying to you''. Even the language barely changes. One side says, ``Unified interpretation maintains order''; the other, ``My conscience is bound by the text I have read''. In today's arguments about vaccines, climate and history, you can almost hear 1517 echo word for word.

Indulgences have not quite disappeared either; the products have changed. Anything for which you pay a price to acquire the reassurance that your account is settled carries their shadow: donations to soothe an uneasy conscience, reposts instead of action, ``expiatory'' purchases to cancel a sense of environmental or interpersonal debt. Thesis 82's question---if this concerns souls, why does the person collecting money not pay?---can still be addressed word for word to many solemn counters.

You possess something our peasant woman did not: the three-strand tool. Next time someone gives a single cause for a great upheaval---``It is all a capitalist conspiracy'', ``Technology determines everything'', ``People have simply become worse''---you know how to unpack it. Ideas, interests and conditions are separate strands. Without any one, the rope cannot pull history.

\section[Your Turn: Where Should Authority Live?]{Your Turn: Where Should Authority Live?}
Return to the paper with which we began.

Five centuries ago, a peasant woman put the money for winter clothes into a chest and received a sheet bearing her husband's name. She believed. Later someone told her: you were deceived; true remission costs nothing; read this book yourself. Later still: before reading, ask who translated it, who chose the words and who paid for its printing. Eventually war taught everyone that when neither side could persuade the other, they needed rules allowing each to manage its own affairs.

Now answer, with reasons: \textbf{where should authority live?}

In institutions? Their unified interpretations have sustained order, and the Church's warning that free interpretation would produce chaos was partly vindicated; today's conspiracy theories offer fresh evidence. In texts? But people must read and translate them. A translator's pen becomes a new sceptre, and ``everyone reading independently'' may mean everyone crowned by their own information bubble. In tested expert consensus? Yet experts once collectively sold indulgences, and the history of science is full of overturned consensuses. In each person's conscience? Luther at Worms commands respect, but when everyone says, ``My conscience says so'', who judges a lawsuit between consciences?

Five hundred years ago, people answered through wars costing millions of lives. Their answer was provisional: each governs its own affairs without interference. Today, when a trending topic can launch an ``everyone do your own research'' movement within hours, is that old answer enough?

There is no standard answer. But notice how you are weighing these arguments: citing evidence, considering counterexamples and fully presenting positions you dislike. That is one genuinely good inheritance from those civil wars.
